%% GENERATED by generate_papers.py -- edit the bank_*.py files, not this file.
\documentclass[11pt,a4paper]{article}
\usepackage[a4paper, top=20mm, bottom=20mm, left=18mm, right=18mm, headheight=14pt]{geometry}
\usepackage[T1]{fontenc}
\usepackage{lmodern}
\usepackage{amsmath, amssymb, mathtools}
\usepackage{braket}
\usepackage{array, tabularx, booktabs, multicol}
\usepackage{enumitem}
\usepackage[dvipsnames]{xcolor}
\usepackage[most]{tcolorbox}
\usepackage{needspace}
\usepackage{fancyhdr, lastpage}
\usepackage{pifont}
\usepackage{tikz}
\usepackage[hidelinks]{hyperref}
\newcolumntype{L}{>{\raggedright\arraybackslash}X}
\definecolor{brand}{HTML}{1F4E79}
\definecolor{okgreen}{HTML}{1B7F3B}
\setlength{\parindent}{0pt}
\setlength{\parskip}{3pt}
\newcommand{\optlabel}[1]{\textbf{(#1)}}
\newcommand{\correctlabel}[1]{\textcolor{okgreen}{\textbf{(#1)}\,\ding{51}}}
\newenvironment{qitem}[1]{\begin{list}{}{\setlength{\leftmargin}{3.3em}\setlength{\labelwidth}{3.1em}%
  \setlength{\labelsep}{0.2em}\setlength{\topsep}{3pt}\setlength{\parsep}{3pt}}\item[\textbf{Q.#1}\hfill]}{\end{list}}
\newcommand{\qsep}{\par\vspace{2pt}{\color{black!15}\hrule}\vspace{4pt}}
\newcommand{\qmeta}[2]{{\footnotesize\color{black!60}\textsc{#1} \;$\cdot$\; #2}\par}
\newtcolorbox{passagebox}[1]{enhanced, breakable, colback=brand!4, colframe=brand!70, fonttitle=\bfseries,
  title={#1}, left=6pt, right=6pt, top=4pt, bottom=4pt}
\newtcolorbox{sectionbanner}{enhanced, colback=brand, colframe=brand, coltext=white, fontupper=\large\bfseries,
  left=8pt, right=8pt, top=6pt, bottom=6pt, sharp corners}
\newtcolorbox{infobox}[1]{enhanced, breakable, colback=black!3, colframe=black!45, fonttitle=\bfseries,
  title={#1}, left=6pt, right=6pt, top=4pt, bottom=4pt}
\newtcolorbox{solbox}[1]{enhanced, breakable, colback=okgreen!3, colframe=okgreen!60!black, fonttitle=\bfseries,
  title={Answer: (#1) \quad Detailed solution}, left=6pt, right=6pt, top=3pt, bottom=4pt, before skip=4pt, after skip=6pt}
\newcommand{\QFTop}{\mathrm{QFT}}
\newcommand{\bub}[1]{\tikz[baseline=-0.65ex]\node[draw, circle, inner sep=0pt, minimum size=4.6mm, font=\scriptsize]{#1};}
\allowdisplaybreaks

\begin{document}
\pagestyle{fancy}\fancyhf{}
\fancyhead[L]{\small\color{black!70}Quantum Fourier Transform --- Assessment}
\fancyhead[R]{\small\color{black!70}\textbf{Question Paper}}
\fancyfoot[C]{\small Page \thepage\ of \pageref*{LastPage}}
\renewcommand{\headrulewidth}{0.4pt}

\thispagestyle{empty}
\begin{center}
{\color{brand}\rule{\linewidth}{1.2pt}}\\[6pt]
{\Huge\bfseries\color{brand} Quantum Fourier Transform}\\[4pt]
{\Large Assessment of the QFT algorithm}\\[4pt]
{\large\bfseries QUESTION PAPER}\\[2pt]
{\color{brand}\rule{\linewidth}{1.2pt}}\\[1pt]
\end{center}
\vspace{-2pt}
\begin{tabularx}{\linewidth}{@{}lXlr@{}}
\textbf{Time allowed:} & 3 hours & \textbf{Maximum marks:} & 240\\
\textbf{Questions:} & 80 (4 sections $\times$ 20) & \textbf{Answer key:} & separate Answer Book
\end{tabularx}
\vspace{6pt}

\textbf{Name of the candidate:} \hrulefill\quad \textbf{Roll No.:} \rule{35mm}{0.4pt}

\vspace{6pt}
\begin{infobox}{General instructions}
\begin{enumerate}[leftmargin=1.4em, itemsep=1pt, topsep=1pt]
\item This question paper has \textbf{four sections} of \textbf{20 questions} each:
\textbf{A}~Mathematics for the QFT, \textbf{B}~Quantum Mechanics, \textbf{C}~QFT Basics, \textbf{D}~Advanced QFT.
\item Every question is of the \textbf{single-correct} type: four options, exactly one correct. Question formats used:
single-correct MCQ, matching lists (List-I/List-II), Assertion--Reason, Statement~I/II and paragraph-based questions.
\item Mark your answers on the \textbf{Response Sheet} at the end of this paper by darkening ONE bubble per question.
\item Calculators are \textbf{not} allowed. Rough work may be done on separate sheets.
\item Unless a question says otherwise, use the \textbf{conventions} printed below; they apply to every section.
\item The detailed solutions are in the separate \emph{Answer Book}; do not consult it before finishing the test.
\end{enumerate}
\end{infobox}
\begin{infobox}{Marking scheme (every question)}
Each question has \textbf{four} options (A), (B), (C) and (D). \textbf{ONLY ONE} of these four options is correct.\\[2pt]
\begin{tabular}{@{}ll@{}}
\emph{Full marks}: & $+3$ if ONLY the correct option is chosen;\\
\emph{Zero marks}: & $\phantom{+}0$ if none of the options is chosen (question unanswered);\\
\emph{Negative marks}: & $-1$ in all other cases (wrong option, or more than one option chosen).
\end{tabular}
\end{infobox}

\begin{infobox}{Conventions and notation used throughout the paper}
\small
\begin{itemize}[leftmargin=1.2em, itemsep=1pt, topsep=2pt]
\item $n$ qubits, $N = 2^n$, $\omega = e^{2\pi i/N}$. A basis state is written $\ket{x} = \ket{x_1x_2\dots x_n}$ with
      $x = \sum_{j} x_j 2^{\,n-j}$: \textbf{qubit 1 is the most significant bit} and is written first.
\item $\displaystyle \mathrm{QFT}\ket{x} = \frac{1}{\sqrt N}\sum_{y=0}^{N-1}\omega^{xy}\ket{y}$, \quad
      $\displaystyle \mathrm{QFT}^{-1}\ket{y} = \frac{1}{\sqrt N}\sum_{x=0}^{N-1}\omega^{-xy}\ket{x}$.
      Binary fraction: $0.x_lx_{l+1}\dots x_m = \frac{x_l}{2} + \frac{x_{l+1}}{4} + \dots + \frac{x_m}{2^{m-l+1}}$.
\item $H = \frac{1}{\sqrt2}\begin{psmallmatrix}1&1\\1&-1\end{psmallmatrix}$; $X, Y, Z$ are the Pauli matrices;
      $R_k = \operatorname{diag}\big(1, e^{2\pi i/2^k}\big)$, so $R_1 = Z$, $S = R_2$, $T = R_3$;
      $\ket{\pm} = \frac{1}{\sqrt2}(\ket0 \pm \ket1)$. Two-qubit matrices use the basis order $\ket{00}, \ket{01}, \ket{10}, \ket{11}$.
\item \textbf{Textbook QFT circuit:} for $j = 1, \dots, n$: apply $H$ to qubit $j$, then for $k = 2, \dots, n-j+1$ a controlled-$R_k$
      with control qubit $j+k-1$ and target qubit $j$; finally SWAP qubits $j$ and $n+1-j$ for $j \le n/2$.
\item \textbf{Phase estimation (QPE):} $U\ket{u} = e^{2\pi i\varphi}\ket{u}$, $0 \le \varphi < 1$; $t$ counting qubits in $\ket{+}$ control
      $U^{2^{t-1}}, \dots, U^{2}, U$; the inverse QFT is applied and the outcome $y$ gives the estimate $\varphi \approx y/2^t$.
\item All measurements are in the computational basis unless stated otherwise; $\log$ means $\log_2$.
\end{itemize}
\end{infobox}
\newpage

\begin{sectionbanner}SECTION A \enspace$\cdot$\enspace Mathematics for the QFT \hfill Q.1--20\end{sectionbanner}

{\small \textbf{Syllabus:} complex numbers, roots of unity, geometric series, the discrete Fourier transform, binary fractions, unitary matrices, continued fractions and modular arithmetic.\\ \textbf{Formats:} Single-correct MCQ (Q.1--14); Matching lists (Q.15--16); Assertion--Reason (Q.17); Statement I / II (Q.18); Paragraph based (Q.19--20). Each question: $+3$ / $0$ / $-1$.}
\par\medskip

\needspace{7\baselineskip}
\begin{qitem}{1}
Let $\omega = e^{2\pi i/8}$. The value of $\displaystyle S = \sum_{k=0}^{7} k\,\omega^{k}$ is
\par\smallskip\noindent\begin{tabularx}{\linewidth}{@{}*{4}{L}@{}}
\optlabel{A}\,$\dfrac{8}{\omega - 1}$ & \optlabel{B}\,$\dfrac{8}{1-\omega}$ & \optlabel{C}\,$0$ & \optlabel{D}\,$\dfrac{8\omega}{\omega-1}$
\end{tabularx}
\end{qitem}
\qsep

\needspace{7\baselineskip}
\begin{qitem}{2}
The number of ordered pairs $(x, y)$ with $x, y \in \{0, 1, \dots, 15\}$ for which
$\displaystyle \sum_{z=0}^{15} e^{2\pi i\,(x+y)z/16} \neq 0$ is
\par\smallskip\noindent\begin{tabularx}{\linewidth}{@{}*{4}{L}@{}}
\optlabel{A}\,$1$ & \optlabel{B}\,$16$ & \optlabel{C}\,$15$ & \optlabel{D}\,$31$
\end{tabularx}
\end{qitem}
\qsep

\needspace{7\baselineskip}
\begin{qitem}{3}
Let $x = 45$, written with $n = 6$ bits as $x_1x_2\dots x_6$. The value of $e^{2\pi i\,x/2^4}$ is
\par\smallskip\noindent\begin{tabularx}{\linewidth}{@{}*{4}{L}@{}}
\optlabel{A}\,$e^{i\,13\pi/8}$ & \optlabel{B}\,$e^{i\,5\pi/8}$ & \optlabel{C}\,$e^{i\,3\pi/8}$ & \optlabel{D}\,$e^{-i\,13\pi/8}$
\end{tabularx}
\end{qitem}
\qsep

\needspace{7\baselineskip}
\begin{qitem}{4}
The matrix $U = \dfrac12\begin{pmatrix}1&1&1&1\\ 1&i&-1&-i\\ 1&-1&1&-1\\ 1&c&-1&d\end{pmatrix}$ is unitary. Then $(c, d)$ equals
\par\smallskip\noindent\begin{tabularx}{\linewidth}{@{}*{4}{L}@{}}
\optlabel{A}\,$(i,\ -i)$ & \optlabel{B}\,$(1,\ -1)$ & \optlabel{C}\,$(-1,\ 1)$ & \optlabel{D}\,$(-i,\ i)$
\end{tabularx}
\end{qitem}
\qsep

\needspace{7\baselineskip}
\begin{qitem}{5}
Let $F$ be the $4\times4$ unitary Fourier matrix with entries $F_{yx} = \tfrac12\, i^{xy}$ ($x, y = 0, \dots, 3$).
The trace of $F$ is
\par\smallskip\noindent\begin{tabularx}{\linewidth}{@{}*{4}{L}@{}}
\optlabel{A}\,$1 + i$ & \optlabel{B}\,$2$ & \optlabel{C}\,$0$ & \optlabel{D}\,$1 - i$
\end{tabularx}
\end{qitem}
\qsep

\needspace{7\baselineskip}
\begin{qitem}{6}
Let $f(\theta) = \big|\,1 + e^{i\theta} + e^{2i\theta} + e^{3i\theta}\,\big|^2$. The maximum value of $f$ and the smallest
positive $\theta$ with $f(\theta) = 0$ are, respectively,
\par\smallskip\noindent\begin{tabularx}{\linewidth}{@{}*{4}{L}@{}}
\optlabel{A}\,$4$ and $\pi/2$ & \optlabel{B}\,$4$ and $\pi$ & \optlabel{C}\,$16$ and $\pi/4$ & \optlabel{D}\,$16$ and $\pi/2$
\end{tabularx}
\end{qitem}
\qsep

\needspace{7\baselineskip}
\begin{qitem}{7}
The best rational approximation $p/q$ to $\dfrac{11}{64}$ with $1 \le q \le 10$ (as produced by the continued-fraction
algorithm) is
\par\smallskip\noindent\begin{tabularx}{\linewidth}{@{}*{4}{L}@{}}
\optlabel{A}\,$2/9$ & \optlabel{B}\,$1/6$ & \optlabel{C}\,$1/7$ & \optlabel{D}\,$1/5$
\end{tabularx}
\end{qitem}
\qsep

\needspace{7\baselineskip}
\begin{qitem}{8}
The order $r$ of $a = 3$ modulo $N = 35$ (the least $r > 0$ with $a^r \equiv 1 \pmod{35}$), and the numbers
$\gcd(a^{r/2} - 1, N)$ and $\gcd(a^{r/2} + 1, N)$, are
\par\smallskip\noindent\begin{tabularx}{\linewidth}{@{}*{2}{L}@{}}
\optlabel{A}\,$r = 6$; $\;5$ and $7$ & \optlabel{B}\,$r = 12$; $\;1$ and $35$ \\[3pt]
\optlabel{C}\,$r = 12$; $\;7$ and $5$ & \optlabel{D}\,$r = 4$; $\;5$ and $7$
\end{tabularx}
\end{qitem}
\qsep

\needspace{7\baselineskip}
\begin{qitem}{9}
A general pure state of 10 qubits is described by $D$ complex amplitudes. After removing normalisation and
the unobservable global phase, the number of independent \emph{real} parameters is $P$. Then $(D, P)$ is
\par\smallskip\noindent\begin{tabularx}{\linewidth}{@{}*{4}{L}@{}}
\optlabel{A}\,$(1024,\ 2048)$ & \optlabel{B}\,$(1024,\ 2046)$ & \optlabel{C}\,$(1024,\ 2047)$ & \optlabel{D}\,$(20,\ 40)$
\end{tabularx}
\end{qitem}
\qsep

\needspace{7\baselineskip}
\begin{qitem}{10}
For $n = 4$, $H^{\otimes 4}|x\rangle = \frac14\sum_{y}(-1)^{x\cdot y}|y\rangle$ where $x\cdot y = \sum_j x_jy_j$.
For $x = 1011_2$, the number of $y \in \{0,\dots,15\}$ with a \emph{negative} amplitude is
\par\smallskip\noindent\begin{tabularx}{\linewidth}{@{}*{4}{L}@{}}
\optlabel{A}\,$4$ & \optlabel{B}\,$8$ & \optlabel{C}\,$12$ & \optlabel{D}\,$3$
\end{tabularx}
\end{qitem}
\qsep

\needspace{7\baselineskip}
\begin{qitem}{11}
A classical FFT on a vector of length $2^n$ needs about $n\,2^n$ operations, while the QFT circuit needs
$\frac{n(n+1)}{2} + \lfloor n/2\rfloor$ gates. For $n = 20$ the ratio (FFT operations)/(QFT gates) is closest to
\par\smallskip\noindent\begin{tabularx}{\linewidth}{@{}*{4}{L}@{}}
\optlabel{A}\,$10^{9}$ & \optlabel{B}\,$10^{3}$ & \optlabel{C}\,$10^{5}$ & \optlabel{D}\,$10^{7}$
\end{tabularx}
\end{qitem}
\qsep

\needspace{7\baselineskip}
\begin{qitem}{12}
With $\omega = e^{2\pi i/8}$, the value of $\displaystyle\prod_{k=1}^{7}\big|1 - \omega^{k}\big|$ is
\par\smallskip\noindent\begin{tabularx}{\linewidth}{@{}*{4}{L}@{}}
\optlabel{A}\,$8$ & \optlabel{B}\,$0$ & \optlabel{C}\,$1$ & \optlabel{D}\,$7$
\end{tabularx}
\end{qitem}
\qsep

\needspace{7\baselineskip}
\begin{qitem}{13}
The normalised vector $\alpha = \frac13(1, 2, 2, 0)$ is transformed by the unitary $4\times4$ Fourier matrix
$\beta_y = \frac12\sum_{x}\alpha_x\, i^{xy}$. The value of $|\beta_0|^2$ is
\par\smallskip\noindent\begin{tabularx}{\linewidth}{@{}*{4}{L}@{}}
\optlabel{A}\,$\dfrac56$ & \optlabel{B}\,$\dfrac{25}{12}$ & \optlabel{C}\,$\dfrac14$ & \optlabel{D}\,$\dfrac{25}{36}$
\end{tabularx}
\end{qitem}
\qsep

\needspace{7\baselineskip}
\begin{qitem}{14}
Let $\varphi = 0.\varphi_1\varphi_2\varphi_3$ (binary, $\varphi_j \in \{0,1\}$). Then $e^{2\pi i\,4\varphi}$ equals
\par\smallskip\noindent\begin{tabularx}{\linewidth}{@{}*{4}{L}@{}}
\optlabel{A}\,$e^{2\pi i\,0.\varphi_2\varphi_3}$ & \optlabel{B}\,$i^{\varphi_2}$ & \optlabel{C}\,$(-1)^{\varphi_1}$ & \optlabel{D}\,$(-1)^{\varphi_3}$
\end{tabularx}
\end{qitem}
\qsep

\needspace{7\baselineskip}
\begin{qitem}{15}
Let $\omega = e^{2\pi i/8}$. Match List-I with List-II.
\par\smallskip\noindent{\renewcommand{\arraystretch}{1.35}%
\begin{tabularx}{\linewidth}{@{}c>{\raggedright\arraybackslash}p{0.33\linewidth}@{\hspace{1.2em}}cL@{}}
\toprule
\multicolumn{2}{@{}l}{\textbf{List-I}} & \multicolumn{2}{l@{}}{\textbf{List-II}} \\ \midrule
\textbf{(P)} & $\omega^{4}$ & \textbf{(1)} & $\sqrt2$ \\
\textbf{(Q)} & $\sum_{k=0}^{7}\omega^{2k}$ & \textbf{(2)} & $-1$ \\
\textbf{(R)} & $\omega\cdot\omega^{7}$ & \textbf{(3)} & $0$ \\
\textbf{(S)} & $\omega + \omega^{7}$ & \textbf{(4)} & $1$ \\
\bottomrule
\end{tabularx}}
\par\smallskip\noindent\begin{tabularx}{\linewidth}{@{}*{2}{L}@{}}
\optlabel{A}\,P$\to$1, Q$\to$3, R$\to$4, S$\to$2 & \optlabel{B}\,P$\to$2, Q$\to$4, R$\to$3, S$\to$1 \\[3pt]
\optlabel{C}\,P$\to$2, Q$\to$3, R$\to$4, S$\to$1 & \optlabel{D}\,P$\to$3, Q$\to$2, R$\to$1, S$\to$4
\end{tabularx}
\end{qitem}
\qsep

\needspace{7\baselineskip}
\begin{qitem}{16}
$F$ is the $4\times4$ unitary Fourier matrix ($F_{yx} = i^{xy}/2$), $H$ the Hadamard matrix and
$CT = \operatorname{diag}(1, 1, 1, e^{i\pi/4})$. Match each matrix in List-I with the \emph{only} property in List-II that it has.
\par\smallskip\noindent{\renewcommand{\arraystretch}{1.35}%
\begin{tabularx}{\linewidth}{@{}c>{\raggedright\arraybackslash}p{0.33\linewidth}@{\hspace{1.2em}}cL@{}}
\toprule
\multicolumn{2}{@{}l}{\textbf{List-I}} & \multicolumn{2}{l@{}}{\textbf{List-II}} \\ \midrule
\textbf{(P)} & $F$ & \textbf{(1)} & a permutation matrix that is its own inverse \\
\textbf{(Q)} & $F^2$ & \textbf{(2)} & real, symmetric, its own inverse, but not a permutation matrix \\
\textbf{(R)} & $H\otimes H$ & \textbf{(3)} & diagonal with determinant $e^{i\pi/4}$ \\
\textbf{(S)} & $CT$ & \textbf{(4)} & $M^4 = I$ but $M^2 \neq I$ \\
\bottomrule
\end{tabularx}}
\par\smallskip\noindent\begin{tabularx}{\linewidth}{@{}*{2}{L}@{}}
\optlabel{A}\,P$\to$4, Q$\to$2, R$\to$1, S$\to$3 & \optlabel{B}\,P$\to$4, Q$\to$1, R$\to$2, S$\to$3 \\[3pt]
\optlabel{C}\,P$\to$3, Q$\to$1, R$\to$2, S$\to$4 & \optlabel{D}\,P$\to$1, Q$\to$4, R$\to$2, S$\to$3
\end{tabularx}
\end{qitem}
\qsep

\needspace{7\baselineskip}
\begin{qitem}{17}
\textbf{Assertion (A):} For every integer $k$ that is not a multiple of $N$, $\displaystyle\sum_{j=0}^{N-1} e^{2\pi i\,jk/N} = 0$.\\[1mm]
\textbf{Reason (R):} For every integer $k$ that is not a multiple of $N$, the $N$ numbers $e^{2\pi i\,jk/N}$
($j = 0, \dots, N-1$) are the $N$ distinct vertices of a regular $N$-gon centred at the origin.
\begin{itemize}[leftmargin=2.6em, labelsep=0.4em, itemsep=2pt, topsep=3pt]
\item[\optlabel{A}] Both Assertion (A) and Reason (R) are true, and R is the correct explanation of A
\item[\optlabel{B}] Both A and R are true, but R is \emph{not} the correct explanation of A
\item[\optlabel{C}] A is true, but R is false
\item[\optlabel{D}] A is false, but R is true
\end{itemize}
\end{qitem}
\qsep

\needspace{7\baselineskip}
\begin{qitem}{18}
\textbf{Statement I:} Every unitary matrix $U$ satisfies $|\det U| = 1$.\\[1mm]
\textbf{Statement II:} The unitary Fourier matrix $F_N$ ($N = 2^n$, entries $\omega^{xy}/\sqrt N$) has $\det F_N = \pm 1$ for every $n$.
\begin{itemize}[leftmargin=2.6em, labelsep=0.4em, itemsep=2pt, topsep=3pt]
\item[\optlabel{A}] Both Statement I and Statement II are true
\item[\optlabel{B}] Both Statement I and Statement II are false
\item[\optlabel{C}] Statement I is true, but Statement II is false
\item[\optlabel{D}] Statement I is false, but Statement II is true
\end{itemize}
\end{qitem}
\qsep

\needspace{20\baselineskip}
\begin{passagebox}{Paragraph for Questions 19 and 20}
Let $N = 32$ and define the sequence $a_x = 1$ if $x \equiv 2 \pmod 5$ and $a_x = 0$ otherwise, for
$x = 0, 1, \dots, 31$. Its (un-normalised) Fourier coefficients are
$A_y = \sum_{x=0}^{31} a_x\, e^{2\pi i\,xy/32}$, $\;y = 0, 1, \dots, 31$. (The sequence is ``periodic with period 5'', but
5 does not divide 32.)
\end{passagebox}
\needspace{7\baselineskip}
\begin{qitem}{19}
The value of $\displaystyle\sum_{y=0}^{31} |A_y|^2$ is
\par\smallskip\noindent\begin{tabularx}{\linewidth}{@{}*{4}{L}@{}}
\optlabel{A}\,$192$ & \optlabel{B}\,$6$ & \optlabel{C}\,$36$ & \optlabel{D}\,$1152$
\end{tabularx}
\end{qitem}
\qsep

\needspace{7\baselineskip}
\begin{qitem}{20}
Among $y = 1, 2, \dots, 15$, the value of $y$ for which $|A_y|$ is largest is
\begin{itemize}[leftmargin=2.6em, labelsep=0.4em, itemsep=2pt, topsep=3pt]
\item[\optlabel{A}] $6$
\item[\optlabel{B}] $16$ is the only peak; all of $1,\dots,15$ give $|A_y| = 0$
\item[\optlabel{C}] $7$
\item[\optlabel{D}] $13$
\end{itemize}
\end{qitem}
\qsep

\newpage

\begin{sectionbanner}SECTION B \enspace$\cdot$\enspace Quantum Mechanics \hfill Q.21--40\end{sectionbanner}

{\small \textbf{Syllabus:} state vectors, Born rule, global and relative phases, gates and unitarity, tensor products, entanglement, partial measurement and phase kickback.\\ \textbf{Formats:} Single-correct MCQ (Q.21--34); Matching lists (Q.35--36); Assertion--Reason (Q.37); Statement I / II (Q.38); Paragraph based (Q.39--40). Each question: $+3$ / $0$ / $-1$.}
\par\medskip

\needspace{7\baselineskip}
\begin{qitem}{21}
A two-qubit state is $\ket{\psi} = c\,\big(\ket{00} + 2i\ket{01} - \ket{10} + \sqrt2\,\ket{11}\big)$ with $c > 0$.
If only the \emph{first} qubit is measured in the computational basis, the probability of obtaining $1$ is
\par\smallskip\noindent\begin{tabularx}{\linewidth}{@{}*{4}{L}@{}}
\optlabel{A}\,$\dfrac12$ & \optlabel{B}\,$\dfrac38$ & \optlabel{C}\,$\dfrac14$ & \optlabel{D}\,$\dfrac58$
\end{tabularx}
\end{qitem}
\qsep

\needspace{7\baselineskip}
\begin{qitem}{22}
Which pair of single-qubit state vectors describes the \emph{same} physical state (i.e.\ gives identical
measurement statistics in every basis)?
\par\smallskip\noindent\begin{tabularx}{\linewidth}{@{}*{2}{L}@{}}
\optlabel{A}\,$\ket{0} + i\ket{1}\;$ and $\;i\ket{0} - \ket{1}$ & \optlabel{B}\,$\ket{0} + i\ket{1}\;$ and $\;i\ket{0} + \ket{1}$ \\[3pt]
\optlabel{C}\,$\ket{0} + i\ket{1}\;$ and $\;\ket{0} - i\ket{1}$ & \optlabel{D}\,$\ket{0} + \ket{1}\;$ and $\;\ket{0} - \ket{1}$
\end{tabularx}
\end{qitem}
\qsep

\needspace{7\baselineskip}
\begin{qitem}{23}
The state $\frac{1}{\sqrt2}\big(\ket{0} + e^{i\phi}\ket{1}\big)$ passes through a Hadamard gate and is then measured
in the computational basis. The probability of outcome $0$ is
\par\smallskip\noindent\begin{tabularx}{\linewidth}{@{}*{4}{L}@{}}
\optlabel{A}\,$\cos^2(\phi/2)$ & \optlabel{B}\,$\sin^2(\phi/2)$ & \optlabel{C}\,$\dfrac12$ & \optlabel{D}\,$\cos^2\phi$
\end{tabularx}
\end{qitem}
\qsep

\needspace{7\baselineskip}
\begin{qitem}{24}
For the qubit state $\ket{\psi} = \cos\frac{\theta}{2}\ket{0} + e^{i\phi}\sin\frac{\theta}{2}\ket{1}$ the quantity
$\langle X\rangle + i\langle Y\rangle$, where $\langle A\rangle = \bra{\psi}A\ket{\psi}$ and $X, Y$ are Pauli matrices, equals
\par\smallskip\noindent\begin{tabularx}{\linewidth}{@{}*{4}{L}@{}}
\optlabel{A}\,$\cos\theta\; e^{i\phi}$ & \optlabel{B}\,$\sin\theta\; e^{i\phi}$ & \optlabel{C}\,$\sin\frac{\theta}{2}\, e^{i\phi}$ & \optlabel{D}\,$\sin\theta\; e^{-i\phi}$
\end{tabularx}
\end{qitem}
\qsep

\needspace{7\baselineskip}
\begin{qitem}{25}
Which of the following two-qubit states is \emph{entangled} (cannot be written as $\ket{a}\otimes\ket{b}$)?
\par\smallskip\noindent\begin{tabularx}{\linewidth}{@{}*{2}{L}@{}}
\optlabel{A}\,$\frac12\big(\ket{00} - \ket{01} + \ket{10} - \ket{11}\big)$ & \optlabel{B}\,$\frac12\big(\ket{00} + \ket{01} + \ket{10} + \ket{11}\big)$ \\[3pt]
\optlabel{C}\,$\frac12\big(\ket{00} + \ket{01} + \ket{10} - \ket{11}\big)$ & \optlabel{D}\,$\frac12\big(\ket{00} + i\ket{01} - \ket{10} - i\ket{11}\big)$
\end{tabularx}
\end{qitem}
\qsep

\needspace{7\baselineskip}
\begin{qitem}{26}
A CNOT gate (qubit 1 = control, qubit 2 = target) acts on $\ket{+}\ket{-}$, where
$\ket{\pm} = \frac{1}{\sqrt2}(\ket0 \pm \ket1)$. The output is
\par\smallskip\noindent\begin{tabularx}{\linewidth}{@{}*{2}{L}@{}}
\optlabel{A}\,$\ket{+}\ket{-}$ & \optlabel{B}\,$\ket{-}\ket{+}$ \\[3pt]
\optlabel{C}\,$\frac{1}{\sqrt2}\big(\ket{0}\ket{-} + \ket{1}\ket{+}\big)$ & \optlabel{D}\,$\ket{-}\ket{-}$
\end{tabularx}
\end{qitem}
\qsep

\needspace{7\baselineskip}
\begin{qitem}{27}
In the basis $\ket{00}, \ket{01}, \ket{10}, \ket{11}$ (first qubit = control), the matrix of controlled-$R_3$, where
$R_k = \operatorname{diag}\big(1, e^{2\pi i/2^k}\big)$, is
\par\smallskip\noindent\begin{tabularx}{\linewidth}{@{}*{2}{L}@{}}
\optlabel{A}\,$\operatorname{diag}\big(1, 1, 1, e^{i\pi/4}\big)$ & \optlabel{B}\,$\operatorname{diag}\big(1, e^{i\pi/4}, 1, e^{i\pi/4}\big)$ \\[3pt]
\optlabel{C}\,$\operatorname{diag}\big(1, 1, e^{i\pi/4}, e^{i\pi/4}\big)$ & \optlabel{D}\,$\operatorname{diag}\big(1, 1, 1, e^{i\pi/8}\big)$
\end{tabularx}
\end{qitem}
\qsep

\needspace{7\baselineskip}
\begin{qitem}{28}
$U\ket{u} = e^{2\pi i\varphi}\ket{u}$. A control qubit in $\ket{+}$ and the register in $\ket{u}$ pass through a
controlled-$U^{2^j}$ gate. The output state is
\par\smallskip\noindent\begin{tabularx}{\linewidth}{@{}*{2}{L}@{}}
\optlabel{A}\,$\frac{1}{\sqrt2}\,e^{2\pi i\,2^j\varphi}\big(\ket0 + \ket1\big)\ket{u}$ & \optlabel{B}\,$\frac{1}{\sqrt2}\big(\ket0 + e^{2\pi i\,2^j\varphi}\ket1\big)\ket{u}$ \\[3pt]
\optlabel{C}\,$\frac{1}{\sqrt2}\big(\ket0 + e^{2\pi i\,j\varphi}\ket1\big)\ket{u}$ & \optlabel{D}\,$\frac{1}{\sqrt2}\big(\ket0 + e^{2\pi i\,\varphi/2^j}\ket1\big)\ket{u}$
\end{tabularx}
\end{qitem}
\qsep

\needspace{7\baselineskip}
\begin{qitem}{29}
The state $\frac{1}{\sqrt3}\big(\ket{00} + \ket{01} + \ket{11}\big)$ is prepared and only the \emph{first} qubit is
measured, giving $0$. The probability of this outcome and the post-measurement state of the second qubit are
\par\smallskip\noindent\begin{tabularx}{\linewidth}{@{}*{4}{L}@{}}
\optlabel{A}\,$\frac23$ and $\ket{+}$ & \optlabel{B}\,$\frac13$ and $\ket{0}$ & \optlabel{C}\,$\frac23$ and $\ket{0}$ & \optlabel{D}\,$\frac13$ and $\ket{+}$
\end{tabularx}
\end{qitem}
\qsep

\needspace{7\baselineskip}
\begin{qitem}{30}
Which of the following matrices can \emph{not} be a quantum gate?
\par\smallskip\noindent\begin{tabularx}{\linewidth}{@{}*{4}{L}@{}}
\optlabel{A}\,$\frac{1}{\sqrt2}\begin{pmatrix}1&1\\ i&i\end{pmatrix}$ & \optlabel{B}\,$\frac{1}{\sqrt2}\begin{pmatrix}1&1\\ 1&-1\end{pmatrix}$ & \optlabel{C}\,$\frac{1}{\sqrt2}\begin{pmatrix}1&i\\ i&1\end{pmatrix}$ & \optlabel{D}\,$\begin{pmatrix}0&1\\ i&0\end{pmatrix}$
\end{tabularx}
\end{qitem}
\qsep

\needspace{7\baselineskip}
\begin{qitem}{31}
Which two-qubit gate does \emph{not} commute with controlled-$R_2 = \operatorname{diag}(1, 1, 1, i)$?
\par\smallskip\noindent\begin{tabularx}{\linewidth}{@{}*{2}{L}@{}}
\optlabel{A}\,$Z\otimes Z$ & \optlabel{B}\,$S\otimes I$, where $S = \operatorname{diag}(1, i)$ \\[3pt]
\optlabel{C}\,$\mathrm{SWAP}$ & \optlabel{D}\,$H\otimes I$
\end{tabularx}
\end{qitem}
\qsep

\needspace{7\baselineskip}
\begin{qitem}{32}
A state-vector simulator stores each amplitude as a single-precision complex number (two 4-byte floats).
The memory needed for the full state of $32$ qubits is
\par\smallskip\noindent\begin{tabularx}{\linewidth}{@{}*{4}{L}@{}}
\optlabel{A}\,$256$ bytes & \optlabel{B}\,$32$ GiB & \optlabel{C}\,$16$ GiB & \optlabel{D}\,$64$ GiB
\end{tabularx}
\end{qitem}
\qsep

\needspace{7\baselineskip}
\begin{qitem}{33}
$H\otimes H$ is applied to the Bell state $\frac{1}{\sqrt2}\big(\ket{00} - \ket{11}\big)$, followed by a measurement of both
qubits. The probability of the outcome $01$ is
\par\smallskip\noindent\begin{tabularx}{\linewidth}{@{}*{4}{L}@{}}
\optlabel{A}\,$0$ & \optlabel{B}\,$\dfrac14$ & \optlabel{C}\,$\dfrac12$ & \optlabel{D}\,$1$
\end{tabularx}
\end{qitem}
\qsep

\needspace{7\baselineskip}
\begin{qitem}{34}
A source emits a qubit that is $\ket0$ with probability $\frac12$ and $\ket1$ with probability $\frac12$ (a classical
coin decides; nobody knows which). A Hadamard gate is applied and the qubit is measured. The probability of outcome $0$ is
\par\smallskip\noindent\begin{tabularx}{\linewidth}{@{}*{4}{L}@{}}
\optlabel{A}\,$1$ & \optlabel{B}\,$\dfrac12$ & \optlabel{C}\,$\dfrac14$ & \optlabel{D}\,$0$
\end{tabularx}
\end{qitem}
\qsep

\needspace{7\baselineskip}
\begin{qitem}{35}
Match each single-qubit gate in List-I with its action in List-II. ($R_k = \operatorname{diag}(1, e^{2\pi i/2^k})$.)
\par\smallskip\noindent{\renewcommand{\arraystretch}{1.35}%
\begin{tabularx}{\linewidth}{@{}c>{\raggedright\arraybackslash}p{0.33\linewidth}@{\hspace{1.2em}}cL@{}}
\toprule
\multicolumn{2}{@{}l}{\textbf{List-I}} & \multicolumn{2}{l@{}}{\textbf{List-II}} \\ \midrule
\textbf{(P)} & $X$ & \textbf{(1)} & $\ket{1} \mapsto \frac{1}{\sqrt2}\big(\ket0 - \ket1\big)$ \\
\textbf{(Q)} & $H$ & \textbf{(2)} & $\ket{1} \mapsto e^{i\pi/4}\ket{1}$ \\
\textbf{(R)} & $R_2$ & \textbf{(3)} & $\ket{0} \leftrightarrow \ket{1}$ \\
\textbf{(S)} & $R_3$ & \textbf{(4)} & $\ket{1} \mapsto i\ket{1}$ \\
\bottomrule
\end{tabularx}}
\par\smallskip\noindent\begin{tabularx}{\linewidth}{@{}*{2}{L}@{}}
\optlabel{A}\,P$\to$3, Q$\to$1, R$\to$2, S$\to$4 & \optlabel{B}\,P$\to$1, Q$\to$3, R$\to$4, S$\to$2 \\[3pt]
\optlabel{C}\,P$\to$3, Q$\to$1, R$\to$4, S$\to$2 & \optlabel{D}\,P$\to$4, Q$\to$1, R$\to$3, S$\to$2
\end{tabularx}
\end{qitem}
\qsep

\needspace{7\baselineskip}
\begin{qitem}{36}
Match each two-qubit state in List-I with its statistics when \emph{both} qubits are measured in the computational
basis (List-II).
\par\smallskip\noindent{\renewcommand{\arraystretch}{1.35}%
\begin{tabularx}{\linewidth}{@{}c>{\raggedright\arraybackslash}p{0.33\linewidth}@{\hspace{1.2em}}cL@{}}
\toprule
\multicolumn{2}{@{}l}{\textbf{List-I}} & \multicolumn{2}{l@{}}{\textbf{List-II}} \\ \midrule
\textbf{(P)} & $\frac{1}{\sqrt2}\big(\ket{00} + \ket{11}\big)$ & \textbf{(1)} & all four outcomes, each with probability $\frac14$ \\
\textbf{(Q)} & $\frac{1}{\sqrt2}\big(\ket{01} - \ket{10}\big)$ & \textbf{(2)} & the two bits are always equal; $00$ and $11$ each $\frac12$ \\
\textbf{(R)} & $\frac12\big(\ket0 + i\ket1\big)\big(\ket0 - \ket1\big)$ & \textbf{(3)} & first bit always $0$; second bit $0$ or $1$ with $\frac12$ each \\
\textbf{(S)} & $\ket{0}\ket{+}$ & \textbf{(4)} & the two bits always differ; $01$ and $10$ each $\frac12$ \\
\bottomrule
\end{tabularx}}
\par\smallskip\noindent\begin{tabularx}{\linewidth}{@{}*{2}{L}@{}}
\optlabel{A}\,P$\to$1, Q$\to$4, R$\to$2, S$\to$3 & \optlabel{B}\,P$\to$4, Q$\to$2, R$\to$1, S$\to$3 \\[3pt]
\optlabel{C}\,P$\to$2, Q$\to$4, R$\to$3, S$\to$1 & \optlabel{D}\,P$\to$2, Q$\to$4, R$\to$1, S$\to$3
\end{tabularx}
\end{qitem}
\qsep

\needspace{7\baselineskip}
\begin{qitem}{37}
\textbf{Assertion (A):} The states $\ket{+}$ and $\ket{-}$ cannot be distinguished by a measurement in the computational
basis $\{\ket0, \ket1\}$.\\[1mm]
\textbf{Reason (R):} $\ket{+}$ and $\ket{-}$ differ only by a global phase.
\begin{itemize}[leftmargin=2.6em, labelsep=0.4em, itemsep=2pt, topsep=3pt]
\item[\optlabel{A}] Both Assertion (A) and Reason (R) are true, and R is the correct explanation of A
\item[\optlabel{B}] Both A and R are true, but R is \emph{not} the correct explanation of A
\item[\optlabel{C}] A is true, but R is false
\item[\optlabel{D}] A is false, but R is true
\end{itemize}
\end{qitem}
\qsep

\needspace{7\baselineskip}
\begin{qitem}{38}
\textbf{Statement I:} The Bell state $\frac{1}{\sqrt2}\big(\ket{00} + \ket{11}\big)$ can be written as $\ket{a}\otimes\ket{b}$ for
suitable single-qubit states $\ket a$, $\ket b$.\\[1mm]
\textbf{Statement II:} A two-qubit state $c_{00}\ket{00} + c_{01}\ket{01} + c_{10}\ket{10} + c_{11}\ket{11}$ is a product state if and only
if $c_{00}c_{11} = c_{01}c_{10}$.
\begin{itemize}[leftmargin=2.6em, labelsep=0.4em, itemsep=2pt, topsep=3pt]
\item[\optlabel{A}] Both Statement I and Statement II are true
\item[\optlabel{B}] Both Statement I and Statement II are false
\item[\optlabel{C}] Statement I is true, but Statement II is false
\item[\optlabel{D}] Statement I is false, but Statement II is true
\end{itemize}
\end{qitem}
\qsep

\needspace{20\baselineskip}
\begin{passagebox}{Paragraph for Questions 39 and 40}
\textbf{Hadamard test.} A control qubit starts in $\ket{0}$ and a target register starts in an eigenstate $\ket{u}$
of a unitary $U$ with $U\ket{u} = e^{2\pi i\varphi}\ket{u}$, $0 \le \varphi < 1$. The circuit applies, in order:
(i) $H$ on the control; (ii) controlled-$U$ (control qubit $\to$ target register); (iii) $H$ on the control;
(iv) a measurement of the control qubit in the computational basis. This is the one-counting-qubit version of quantum
phase estimation.
\end{passagebox}
\needspace{7\baselineskip}
\begin{qitem}{39}
The probability that the control qubit is found in state $\ket1$ is
\par\smallskip\noindent\begin{tabularx}{\linewidth}{@{}*{4}{L}@{}}
\optlabel{A}\,$\cos^2(\pi\varphi)$ & \optlabel{B}\,$\sin^2(\pi\varphi/2)$ & \optlabel{C}\,$\sin^2(2\pi\varphi)$ & \optlabel{D}\,$\sin^2(\pi\varphi)$
\end{tabularx}
\end{qitem}
\qsep

\needspace{7\baselineskip}
\begin{qitem}{40}
Take $\varphi = \frac13$. The circuit is run repeatedly (fresh $\ket0$ control each time). The minimum number of runs
needed so that the probability of observing outcome $1$ \emph{at least once} exceeds $0.99$ is
\par\smallskip\noindent\begin{tabularx}{\linewidth}{@{}*{4}{L}@{}}
\optlabel{A}\,$2$ & \optlabel{B}\,$3$ & \optlabel{C}\,$4$ & \optlabel{D}\,$7$
\end{tabularx}
\end{qitem}
\qsep

\newpage

\begin{sectionbanner}SECTION C \enspace$\cdot$\enspace QFT Basics \hfill Q.41--60\end{sectionbanner}

{\small \textbf{Syllabus:} definition and matrix of the QFT, product formula, textbook circuit, gate counts, inverse QFT and simple inputs.\\ \textbf{Formats:} Single-correct MCQ (Q.41--54); Matching lists (Q.55--56); Assertion--Reason (Q.57); Statement I / II (Q.58); Paragraph based (Q.59--60). Each question: $+3$ / $0$ / $-1$.}
\par\medskip

\needspace{7\baselineskip}
\begin{qitem}{41}
For $n = 2$, the amplitude of $\ket{10}$ in $\mathrm{QFT}\ket{11}$ is
\par\smallskip\noindent\begin{tabularx}{\linewidth}{@{}*{4}{L}@{}}
\optlabel{A}\,$-\dfrac12$ & \optlabel{B}\,$\dfrac12$ & \optlabel{C}\,$-\dfrac{i}{2}$ & \optlabel{D}\,$\dfrac{i}{2}$
\end{tabularx}
\end{qitem}
\qsep

\needspace{7\baselineskip}
\begin{qitem}{42}
A 5-qubit QFT is applied to the basis state $\ket{x}$ with $x = 19$ and the output is measured. The probability of
obtaining $y = 7$ is
\par\smallskip\noindent\begin{tabularx}{\linewidth}{@{}*{4}{L}@{}}
\optlabel{A}\,$\dfrac{7}{32}$ & \optlabel{B}\,$\dfrac{1}{\sqrt{32}}$ & \optlabel{C}\,$0$ & \optlabel{D}\,$\dfrac{1}{32}$
\end{tabularx}
\end{qitem}
\qsep

\needspace{7\baselineskip}
\begin{qitem}{43}
For $n = 3$, $\mathrm{QFT}\ket{110}$ equals (qubit 1 written first)
\par\smallskip\noindent\begin{tabularx}{\linewidth}{@{}*{2}{L}@{}}
\optlabel{A}\,$\frac{1}{\sqrt8}\big(\ket0 - \ket1\big)\big(\ket0 - \ket1\big)\big(\ket0 - i\ket1\big)$ & \optlabel{B}\,$\frac{1}{\sqrt8}\big(\ket0 + \ket1\big)\big(\ket0 - \ket1\big)\big(\ket0 + i\ket1\big)$ \\[3pt]
\optlabel{C}\,$\frac{1}{\sqrt8}\big(\ket0 + \ket1\big)\big(\ket0 - \ket1\big)\big(\ket0 - i\ket1\big)$ & \optlabel{D}\,$\frac{1}{\sqrt8}\big(\ket0 - i\ket1\big)\big(\ket0 - \ket1\big)\big(\ket0 + \ket1\big)$
\end{tabularx}
\end{qitem}
\qsep

\needspace{7\baselineskip}
\begin{qitem}{44}
The textbook QFT circuit on $n = 12$ qubits contains how many controlled-$R_k$ gates?
\par\smallskip\noindent\begin{tabularx}{\linewidth}{@{}*{4}{L}@{}}
\optlabel{A}\,$144$ & \optlabel{B}\,$72$ & \optlabel{C}\,$66$ & \optlabel{D}\,$78$
\end{tabularx}
\end{qitem}
\qsep

\needspace{7\baselineskip}
\begin{qitem}{45}
On hardware a SWAP is implemented with 3 CNOTs. Counting every controlled-$R_k$ as one two-qubit gate, the total number
of two-qubit gates in the 12-qubit textbook QFT (including the final SWAPs) is
\par\smallskip\noindent\begin{tabularx}{\linewidth}{@{}*{4}{L}@{}}
\optlabel{A}\,$72$ & \optlabel{B}\,$84$ & \optlabel{C}\,$102$ & \optlabel{D}\,$66$
\end{tabularx}
\end{qitem}
\qsep

\needspace{7\baselineskip}
\begin{qitem}{46}
The smallest rotation angle appearing in the textbook QFT circuit on $6$ qubits is
\par\smallskip\noindent\begin{tabularx}{\linewidth}{@{}*{4}{L}@{}}
\optlabel{A}\,$11.25^\circ$ & \optlabel{B}\,$6^\circ$ & \optlabel{C}\,$2.8125^\circ$ & \optlabel{D}\,$5.625^\circ$
\end{tabularx}
\end{qitem}
\qsep

\needspace{7\baselineskip}
\begin{qitem}{47}
In the textbook QFT circuit acting on $\ket{x_1x_2\dots x_n}$, \emph{before} the final SWAPs, qubit 1 is in the state
$\frac{1}{\sqrt2}\big(\ket0 + e^{2\pi i\,\phi}\ket1\big)$ with $\phi$ equal to
\par\smallskip\noindent\begin{tabularx}{\linewidth}{@{}*{4}{L}@{}}
\optlabel{A}\,$0.x_1x_2\dots x_n$ & \optlabel{B}\,$0.x_nx_{n-1}\dots x_1$ & \optlabel{C}\,$0.x_n$ & \optlabel{D}\,$0.x_1$
\end{tabularx}
\end{qitem}
\qsep

\needspace{7\baselineskip}
\begin{qitem}{48}
For any $n$-qubit basis state, $\mathrm{QFT}^2\ket{x}$ equals ($N = 2^n$)
\par\smallskip\noindent\begin{tabularx}{\linewidth}{@{}*{4}{L}@{}}
\optlabel{A}\,$\ket{x}$ & \optlabel{B}\,$\ket{2x \bmod N}$ & \optlabel{C}\,$\ket{(-x) \bmod N}$ & \optlabel{D}\,$\ket{N - 1 - x}$
\end{tabularx}
\end{qitem}
\qsep

\needspace{7\baselineskip}
\begin{qitem}{49}
The QFT is applied to the uniform superposition $\frac{1}{\sqrt N}\sum_{x=0}^{N-1}\ket{x}$. The output is
\par\smallskip\noindent\begin{tabularx}{\linewidth}{@{}*{2}{L}@{}}
\optlabel{A}\,$\ket{N-1}$ & \optlabel{B}\,$\ket{0}$ \\[3pt]
\optlabel{C}\,the uniform superposition itself & \optlabel{D}\,$\ket{N/2}$
\end{tabularx}
\end{qitem}
\qsep

\needspace{7\baselineskip}
\begin{qitem}{50}
For $n$ qubits, $\mathrm{QFT}\ket{10\dots0}$ (i.e.\ $x = 2^{n-1}$) equals
\par\smallskip\noindent\begin{tabularx}{\linewidth}{@{}*{4}{L}@{}}
\optlabel{A}\,$\ket{+}^{\otimes(n-1)}\otimes\ket{-}$ & \optlabel{B}\,$\ket{-}^{\otimes n}$ & \optlabel{C}\,$\ket{+}^{\otimes n}$ & \optlabel{D}\,$\ket{-}\otimes\ket{+}^{\otimes(n-1)}$
\end{tabularx}
\end{qitem}
\qsep

\needspace{7\baselineskip}
\begin{qitem}{51}
The inverse QFT on 3 qubits is built by reversing the textbook QFT circuit and conjugating every gate. Ignoring the
SWAP, the \emph{first} gate of the inverse circuit is
\par\smallskip\noindent\begin{tabularx}{\linewidth}{@{}*{2}{L}@{}}
\optlabel{A}\,$H$ on qubit 3 & \optlabel{B}\,controlled-$R_3^\dagger$ between qubits 1 and 3 \\[3pt]
\optlabel{C}\,controlled-$R_2^\dagger$ between qubits 1 and 2 & \optlabel{D}\,$H$ on qubit 1
\end{tabularx}
\end{qitem}
\qsep

\needspace{7\baselineskip}
\begin{qitem}{52}
For $n = 4$, the phase (in degrees, in $[0^\circ, 360^\circ)$) of the amplitude of $\ket{y}$, $y = 3$, in $\mathrm{QFT}\ket{x}$, $x = 5$, is
\par\smallskip\noindent\begin{tabularx}{\linewidth}{@{}*{4}{L}@{}}
\optlabel{A}\,$22.5^\circ$ & \optlabel{B}\,$337.5^\circ$ & \optlabel{C}\,$168.75^\circ$ & \optlabel{D}\,$15^\circ$
\end{tabularx}
\end{qitem}
\qsep

\needspace{7\baselineskip}
\begin{qitem}{53}
NumPy's two transforms of a length-$N$ vector $a$ are
\[ \texttt{fft}(a)_y = \sum_x a_x\, e^{-2\pi i\,xy/N}, \qquad \texttt{ifft}(a)_y = \frac1N\sum_x a_x\, e^{+2\pi i\,xy/N}. \]
If $\mathrm{QFT}\sum_x\alpha_x\ket x = \sum_y\beta_y\ket y$ (convention box), then the vector $\beta$ equals
\par\smallskip\noindent\begin{tabularx}{\linewidth}{@{}*{4}{L}@{}}
\optlabel{A}\,$\texttt{fft}(\alpha)$ & \optlabel{B}\,$\sqrt N\cdot\texttt{ifft}(\alpha)$ & \optlabel{C}\,$\texttt{ifft}(\alpha)/\sqrt N$ & \optlabel{D}\,$\texttt{fft}(\alpha)/\sqrt N$
\end{tabularx}
\end{qitem}
\qsep

\needspace{7\baselineskip}
\begin{qitem}{54}
For an $n$-qubit computational basis state $\ket x$, the state $\mathrm{QFT}\ket x$ is
\par\smallskip\noindent\begin{tabularx}{\linewidth}{@{}*{2}{L}@{}}
\optlabel{A}\,a product state only when $x$ is even & \optlabel{B}\,always maximally entangled \\[3pt]
\optlabel{C}\,entangled unless $x = 0$ & \optlabel{D}\,always a product (unentangled) state
\end{tabularx}
\end{qitem}
\qsep

\needspace{7\baselineskip}
\begin{qitem}{55}
Match the number of qubits $n$ (List-I) with the number of controlled-$R_k$ gates in the textbook QFT (List-II).
\par\smallskip\noindent{\renewcommand{\arraystretch}{1.35}%
\begin{tabularx}{\linewidth}{@{}c>{\raggedright\arraybackslash}p{0.33\linewidth}@{\hspace{1.2em}}cL@{}}
\toprule
\multicolumn{2}{@{}l}{\textbf{List-I}} & \multicolumn{2}{l@{}}{\textbf{List-II}} \\ \midrule
\textbf{(P)} & $n = 4$ & \textbf{(1)} & $15$ \\
\textbf{(Q)} & $n = 6$ & \textbf{(2)} & $28$ \\
\textbf{(R)} & $n = 8$ & \textbf{(3)} & $45$ \\
\textbf{(S)} & $n = 10$ & \textbf{(4)} & $6$ \\
\bottomrule
\end{tabularx}}
\par\smallskip\noindent\begin{tabularx}{\linewidth}{@{}*{2}{L}@{}}
\optlabel{A}\,P$\to$1, Q$\to$4, R$\to$2, S$\to$3 & \optlabel{B}\,P$\to$4, Q$\to$2, R$\to$1, S$\to$3 \\[3pt]
\optlabel{C}\,P$\to$4, Q$\to$1, R$\to$2, S$\to$3 & \optlabel{D}\,P$\to$4, Q$\to$1, R$\to$3, S$\to$2
\end{tabularx}
\end{qitem}
\qsep

\needspace{7\baselineskip}
\begin{qitem}{56}
For $n = 2$, match each input (List-I) with $\mathrm{QFT}$ of it (List-II); qubit 1 is written first.
\par\smallskip\noindent{\renewcommand{\arraystretch}{1.35}%
\begin{tabularx}{\linewidth}{@{}c>{\raggedright\arraybackslash}p{0.33\linewidth}@{\hspace{1.2em}}cL@{}}
\toprule
\multicolumn{2}{@{}l}{\textbf{List-I}} & \multicolumn{2}{l@{}}{\textbf{List-II}} \\ \midrule
\textbf{(P)} & $\ket{00}$ & \textbf{(1)} & $\frac12\big(\ket0 - \ket1\big)\big(\ket0 + i\ket1\big)$ \\
\textbf{(Q)} & $\ket{01}$ & \textbf{(2)} & $\frac12\big(\ket0 - \ket1\big)\big(\ket0 - i\ket1\big)$ \\
\textbf{(R)} & $\ket{10}$ & \textbf{(3)} & $\frac12\big(\ket0 + \ket1\big)\big(\ket0 + \ket1\big)$ \\
\textbf{(S)} & $\ket{11}$ & \textbf{(4)} & $\frac12\big(\ket0 + \ket1\big)\big(\ket0 - \ket1\big)$ \\
\bottomrule
\end{tabularx}}
\par\smallskip\noindent\begin{tabularx}{\linewidth}{@{}*{2}{L}@{}}
\optlabel{A}\,P$\to$3, Q$\to$4, R$\to$1, S$\to$2 & \optlabel{B}\,P$\to$4, Q$\to$1, R$\to$3, S$\to$2 \\[3pt]
\optlabel{C}\,P$\to$3, Q$\to$2, R$\to$4, S$\to$1 & \optlabel{D}\,P$\to$3, Q$\to$1, R$\to$4, S$\to$2
\end{tabularx}
\end{qitem}
\qsep

\needspace{7\baselineskip}
\begin{qitem}{57}
\textbf{Assertion (A):} Measuring $\mathrm{QFT}\ket{x}$ in the computational basis reveals no information about the
basis state $x$.\\[1mm]
\textbf{Reason (R):} All $N$ amplitudes of $\mathrm{QFT}\ket{x}$ have the same magnitude $1/\sqrt N$.
\begin{itemize}[leftmargin=2.6em, labelsep=0.4em, itemsep=2pt, topsep=3pt]
\item[\optlabel{A}] Both Assertion (A) and Reason (R) are true, and R is the correct explanation of A
\item[\optlabel{B}] Both A and R are true, but R is \emph{not} the correct explanation of A
\item[\optlabel{C}] A is true, but R is false
\item[\optlabel{D}] A is false, but R is true
\end{itemize}
\end{qitem}
\qsep

\needspace{7\baselineskip}
\begin{qitem}{58}
\textbf{Statement I:} The QFT on $n$ qubits can be implemented with $O(n^2)$ one- and two-qubit gates.\\[1mm]
\textbf{Statement II:} Consequently, for an arbitrary classical input vector of length $2^n$, a quantum computer can output all
$2^n$ of its Fourier coefficients using only $O(n^2)$ operations in total.
\begin{itemize}[leftmargin=2.6em, labelsep=0.4em, itemsep=2pt, topsep=3pt]
\item[\optlabel{A}] Both Statement I and Statement II are true
\item[\optlabel{B}] Both Statement I and Statement II are false
\item[\optlabel{C}] Statement I is true, but Statement II is false
\item[\optlabel{D}] Statement I is false, but Statement II is true
\end{itemize}
\end{qitem}
\qsep

\needspace{20\baselineskip}
\begin{passagebox}{Paragraph for Questions 59 and 60}
The textbook 3-qubit QFT circuit is applied to the input $\ket{x_1x_2x_3} = \ket{011}$ ($x = 3$). The gates, in order, are:
$H$ on qubit 1; controlled-$R_2$ (control qubit 2, target qubit 1); controlled-$R_3$ (control qubit 3, target qubit 1);
$H$ on qubit 2; controlled-$R_2$ (control qubit 3, target qubit 2); $H$ on qubit 3; SWAP of qubits 1 and 3.
Here $R_k = \operatorname{diag}(1, e^{2\pi i/2^k})$, and ``the phase of qubit $j$'' means the angle $\theta$ in its state
$\frac{1}{\sqrt2}(\ket0 + e^{i\theta}\ket1)$, taken in $[0^\circ, 360^\circ)$.
\end{passagebox}
\needspace{7\baselineskip}
\begin{qitem}{59}
Immediately after the controlled-$R_3$ gate (the third gate), the phase of qubit 1 is
\par\smallskip\noindent\begin{tabularx}{\linewidth}{@{}*{4}{L}@{}}
\optlabel{A}\,$270^\circ$ & \optlabel{B}\,$135^\circ$ & \optlabel{C}\,$90^\circ$ & \optlabel{D}\,$45^\circ$
\end{tabularx}
\end{qitem}
\qsep

\needspace{7\baselineskip}
\begin{qitem}{60}
After the complete circuit (including the SWAP), the amplitude of the basis state $\ket{110}$ is
\par\smallskip\noindent\begin{tabularx}{\linewidth}{@{}*{4}{L}@{}}
\optlabel{A}\,$\dfrac{i}{8}$ & \optlabel{B}\,$\dfrac{e^{i\pi/4}}{\sqrt8}$ & \optlabel{C}\,$\dfrac{-i}{\sqrt8}$ & \optlabel{D}\,$\dfrac{i}{\sqrt8}$
\end{tabularx}
\end{qitem}
\qsep

\newpage

\begin{sectionbanner}SECTION D \enspace$\cdot$\enspace Advanced QFT \hfill Q.61--80\end{sectionbanner}

{\small \textbf{Syllabus:} periodic states, shift theorem, phase estimation accuracy, Shor's algorithm, approximate and semiclassical QFT, spectral properties and applications.\\ \textbf{Formats:} Single-correct MCQ (Q.61--74); Matching lists (Q.75--76); Assertion--Reason (Q.77); Statement I / II (Q.78); Paragraph based (Q.79--80). Each question: $+3$ / $0$ / $-1$.}
\par\medskip

\needspace{7\baselineskip}
\begin{qitem}{61}
On $n = 6$ qubits ($N = 64$) the state $\frac{1}{\sqrt8}\sum_{j=0}^{7}\ket{3 + 8j}$ is prepared, the QFT is applied and all
qubits are measured. The possible outcomes $y$ and their probabilities are
\begin{itemize}[leftmargin=2.6em, labelsep=0.4em, itemsep=2pt, topsep=3pt]
\item[\optlabel{A}] only $y = 8$, with probability $1$
\item[\optlabel{B}] all $64$ values of $y$, each with probability $\frac{1}{64}$
\item[\optlabel{C}] $y \in \{3, 11, 19, \dots, 59\}$, each with probability $\frac18$
\item[\optlabel{D}] $y \in \{0, 8, 16, \dots, 56\}$, each with probability $\frac18$
\end{itemize}
\end{qitem}
\qsep

\needspace{7\baselineskip}
\begin{qitem}{62}
For the state $\frac{1}{\sqrt8}\sum_{j=0}^{7}\ket{3 + 8j}$ on $6$ qubits, the phase (in $[0^\circ, 360^\circ)$) of the amplitude of
$\ket{16}$ after the QFT is
\par\smallskip\noindent\begin{tabularx}{\linewidth}{@{}*{4}{L}@{}}
\optlabel{A}\,$90^\circ$ & \optlabel{B}\,$270^\circ$ & \optlabel{C}\,$135^\circ$ & \optlabel{D}\,$0^\circ$
\end{tabularx}
\end{qitem}
\qsep

\needspace{7\baselineskip}
\begin{qitem}{63}
Let $\mathrm{QFT}\sum_x\alpha_x\ket{x} = \sum_y\beta_y\ket{y}$ on $n$ qubits ($N = 2^n$, $\omega = e^{2\pi i/N}$). For a fixed integer $s$,
$\mathrm{QFT}\sum_x\alpha_x\ket{(x + s) \bmod N}$ equals
\par\smallskip\noindent\begin{tabularx}{\linewidth}{@{}*{2}{L}@{}}
\optlabel{A}\,$\sum_y\omega^{-sy}\beta_y\ket{y}$ & \optlabel{B}\,$\sum_y\beta_{(y+s)\bmod N}\ket{y}$ \\[3pt]
\optlabel{C}\,$\sum_y\beta_y\ket{(y + s)\bmod N}$ & \optlabel{D}\,$\sum_y\omega^{sy}\beta_y\ket{y}$
\end{tabularx}
\end{qitem}
\qsep

\needspace{7\baselineskip}
\begin{qitem}{64}
To obtain $\varphi$ accurate to $n$ bits with probability at least $1 - \varepsilon$, phase estimation uses
$t = n + \big\lceil \log_2\!\big(2 + \frac{1}{2\varepsilon}\big)\big\rceil$ counting qubits. For $n = 5$ and $\varepsilon = 0.1$, $t$ equals
\par\smallskip\noindent\begin{tabularx}{\linewidth}{@{}*{4}{L}@{}}
\optlabel{A}\,$7$ & \optlabel{B}\,$8$ & \optlabel{C}\,$10$ & \optlabel{D}\,$9$
\end{tabularx}
\end{qitem}
\qsep

\needspace{7\baselineskip}
\begin{qitem}{65}
Phase estimation with $t = 3$ counting qubits (inverse QFT readout) is run on an eigenstate with $\varphi = \frac58$.
The probability that the counting register reads $101$ is
\par\smallskip\noindent\begin{tabularx}{\linewidth}{@{}*{4}{L}@{}}
\optlabel{A}\,$\dfrac58$ & \optlabel{B}\,$1$ & \optlabel{C}\,$\dfrac18$ & \optlabel{D}\,$\dfrac{4}{\pi^2}$
\end{tabularx}
\end{qitem}
\qsep

\needspace{7\baselineskip}
\begin{qitem}{66}
If $\varphi$ is not exactly a $t$-bit binary fraction, phase estimation returns the \emph{single} closest $t$-bit
approximation with probability at least
\par\smallskip\noindent\begin{tabularx}{\linewidth}{@{}*{4}{L}@{}}
\optlabel{A}\,$\dfrac12$ & \optlabel{B}\,$\dfrac1\pi \approx 0.318$ & \optlabel{C}\,$\dfrac{4}{\pi^2} \approx 0.405$ & \optlabel{D}\,$\dfrac{8}{\pi^2} \approx 0.81$
\end{tabularx}
\end{qitem}
\qsep

\needspace{7\baselineskip}
\begin{qitem}{67}
Phase estimation with $t = 2$ counting qubits is run with $\varphi = \frac13$. The probability of the outcome $y = 1$
(i.e.\ the estimate $\varphi \approx \frac14$) is
\par\smallskip\noindent\begin{tabularx}{\linewidth}{@{}*{4}{L}@{}}
\optlabel{A}\,$\dfrac{6 - 3\sqrt3}{16}$ & \optlabel{B}\,$\dfrac13$ & \optlabel{C}\,$\dfrac{4}{\pi^2}$ & \optlabel{D}\,$\dfrac{6 + 3\sqrt3}{16}$
\end{tabularx}
\end{qitem}
\qsep

\needspace{7\baselineskip}
\begin{qitem}{68}
For Shor's algorithm with a work register of $n = \lceil\log_2 N\rceil$ qubits and a counting register of $2n$ qubits,
the total number of qubits needed to factor $N = 21$ is
\par\smallskip\noindent\begin{tabularx}{\linewidth}{@{}*{4}{L}@{}}
\optlabel{A}\,$15$ & \optlabel{B}\,$12$ & \optlabel{C}\,$21$ & \optlabel{D}\,$10$
\end{tabularx}
\end{qitem}
\qsep

\needspace{7\baselineskip}
\begin{qitem}{69}
Order finding for $a = 7$, $N = 15$ uses a counting register of $t = 8$ qubits ($M = 256$). The set of possible
measurement outcomes $y$ is
\par\smallskip\noindent\begin{tabularx}{\linewidth}{@{}*{4}{L}@{}}
\optlabel{A}\,$\{0, 128\}$ & \optlabel{B}\,$\{0, 64, 128, 192\}$ & \optlabel{C}\,$\{0, 4, 8, 12\}$ & \optlabel{D}\,$\{0, 85, 171\}$
\end{tabularx}
\end{qitem}
\qsep

\needspace{7\baselineskip}
\begin{qitem}{70}
In the setting of the previous question (order finding for $a = 7$, $N = 15$, $M = 256$), the continued-fraction step
turns $y/256$ into $k/r'$. Which outcome leads to the \emph{wrong} candidate $r' = 2$?
\par\smallskip\noindent\begin{tabularx}{\linewidth}{@{}*{4}{L}@{}}
\optlabel{A}\,$y = 64$ & \optlabel{B}\,$y = 128$ & \optlabel{C}\,$y = 0$ & \optlabel{D}\,$y = 192$
\end{tabularx}
\end{qitem}
\qsep

\needspace{7\baselineskip}
\begin{qitem}{71}
An approximate QFT on $n = 20$ qubits keeps only the controlled-$R_k$ gates with $k \le m = 6$ and drops the rest.
The number of controlled rotations dropped is
\par\smallskip\noindent\begin{tabularx}{\linewidth}{@{}*{4}{L}@{}}
\optlabel{A}\,$114$ & \optlabel{B}\,$85$ & \optlabel{C}\,$95$ & \optlabel{D}\,$105$
\end{tabularx}
\end{qitem}
\qsep

\needspace{7\baselineskip}
\begin{qitem}{72}
On $n = 6$ qubits ($N = 64$) the uniform superposition of all $x \equiv 0 \pmod 6$, $0 \le x \le 63$, is prepared and the
QFT applied. The probability of the outcome $y = 32$ is
\par\smallskip\noindent\begin{tabularx}{\linewidth}{@{}*{4}{L}@{}}
\optlabel{A}\,$\dfrac{11}{64}$ & \optlabel{B}\,$\dfrac{1}{64}$ & \optlabel{C}\,$\dfrac16$ & \optlabel{D}\,$\dfrac{1}{11}$
\end{tabularx}
\end{qitem}
\qsep

\needspace{7\baselineskip}
\begin{qitem}{73}
Every eigenvalue $\lambda$ of the $n$-qubit QFT matrix ($n \ge 1$) satisfies
\begin{itemize}[leftmargin=2.6em, labelsep=0.4em, itemsep=2pt, topsep=3pt]
\item[\optlabel{A}] $\lambda^{N} = 1$, and all $N$ of the $N$-th roots of unity occur
\item[\optlabel{B}] $|\lambda| < 1$ for at least one eigenvalue
\item[\optlabel{C}] $\lambda \in \{1, -1\}$
\item[\optlabel{D}] $\lambda \in \{1, -1, i, -i\}$
\end{itemize}
\end{qitem}
\qsep

\needspace{7\baselineskip}
\begin{qitem}{74}
In phase estimation the inverse QFT is immediately followed by measurement of every counting qubit. Using the
semiclassical (Griffiths--Niu) QFT, the number of \emph{two-qubit} gates needed for this inverse QFT on $n$ qubits is
\par\smallskip\noindent\begin{tabularx}{\linewidth}{@{}*{4}{L}@{}}
\optlabel{A}\,$n - 1$ & \optlabel{B}\,$\dfrac{n(n-1)}{2}$ & \optlabel{C}\,$0$ & \optlabel{D}\,$n$
\end{tabularx}
\end{qitem}
\qsep

\needspace{7\baselineskip}
\begin{qitem}{75}
Match each algorithm in List-I with the role played by the QFT, or the related step, in List-II.
\par\smallskip\noindent{\renewcommand{\arraystretch}{1.35}%
\begin{tabularx}{\linewidth}{@{}c>{\raggedright\arraybackslash}p{0.33\linewidth}@{\hspace{1.2em}}cL@{}}
\toprule
\multicolumn{2}{@{}l}{\textbf{List-I}} & \multicolumn{2}{l@{}}{\textbf{List-II}} \\ \midrule
\textbf{(P)} & Quantum phase estimation & \textbf{(1)} & adds a number by applying controlled phase rotations to a register held in the Fourier basis \\
\textbf{(Q)} & Shor's factoring algorithm & \textbf{(2)} & turns the phases $e^{2\pi i\,2^j\varphi}$ in the counting register into the binary digits of $\varphi$ \\
\textbf{(R)} & Draper's quantum adder & \textbf{(3)} & after a QPE step, applies an ancilla rotation proportional to $1/\lambda$ for each eigenvalue $\lambda$ \\
\textbf{(S)} & HHL linear-system solver & \textbf{(4)} & extracts an order $r$ from peaks near $kM/r$ using continued fractions \\
\bottomrule
\end{tabularx}}
\par\smallskip\noindent\begin{tabularx}{\linewidth}{@{}*{2}{L}@{}}
\optlabel{A}\,P$\to$2, Q$\to$4, R$\to$1, S$\to$3 & \optlabel{B}\,P$\to$4, Q$\to$2, R$\to$1, S$\to$3 \\[3pt]
\optlabel{C}\,P$\to$2, Q$\to$1, R$\to$4, S$\to$3 & \optlabel{D}\,P$\to$2, Q$\to$4, R$\to$3, S$\to$1
\end{tabularx}
\end{qitem}
\qsep

\needspace{7\baselineskip}
\begin{qitem}{76}
On $n = 4$ qubits ($N = 16$) the QFT is applied to each input in List-I and the result measured. Match with List-II.
\par\smallskip\noindent{\renewcommand{\arraystretch}{1.35}%
\begin{tabularx}{\linewidth}{@{}c>{\raggedright\arraybackslash}p{0.33\linewidth}@{\hspace{1.2em}}cL@{}}
\toprule
\multicolumn{2}{@{}l}{\textbf{List-I}} & \multicolumn{2}{l@{}}{\textbf{List-II}} \\ \midrule
\textbf{(P)} & $\ket{0}$ & \textbf{(1)} & only $y = 0$ and $y = 8$, each with probability $\frac12$ \\
\textbf{(Q)} & $\frac12\big(\ket0 + \ket4 + \ket8 + \ket{12}\big)$ & \textbf{(2)} & all 16 outcomes, each with probability $\frac{1}{16}$ \\
\textbf{(R)} & $\frac{1}{\sqrt8}\sum_{x\ \mathrm{even}}\ket{x}$ & \textbf{(3)} & only $y = 0, 4, 8, 12$, each with probability $\frac14$ \\
\textbf{(S)} & $\frac14\sum_{x=0}^{15}\ket{x}$ & \textbf{(4)} & $y = 0$ with certainty \\
\bottomrule
\end{tabularx}}
\par\smallskip\noindent\begin{tabularx}{\linewidth}{@{}*{2}{L}@{}}
\optlabel{A}\,P$\to$4, Q$\to$3, R$\to$1, S$\to$2 & \optlabel{B}\,P$\to$2, Q$\to$3, R$\to$4, S$\to$1 \\[3pt]
\optlabel{C}\,P$\to$2, Q$\to$3, R$\to$1, S$\to$4 & \optlabel{D}\,P$\to$2, Q$\to$1, R$\to$3, S$\to$4
\end{tabularx}
\end{qitem}
\qsep

\needspace{7\baselineskip}
\begin{qitem}{77}
\textbf{Assertion (A):} For large $n$, dropping every controlled-$R_k$ with $k > m$, where $m$ is only slightly larger than
$\log_2 n$, changes the output of the QFT circuit only slightly.\\[1mm]
\textbf{Reason (R):} Each dropped gate differs from the identity only by a phase $2\pi/2^{k}$, so its error is at most about
$2\pi/2^{k}$; at most $n$ gates are dropped for each $k$, giving a total error of order $n\cdot 2\pi/2^{m}$, which is small once
$2^m \gg n$.
\begin{itemize}[leftmargin=2.6em, labelsep=0.4em, itemsep=2pt, topsep=3pt]
\item[\optlabel{A}] Both Assertion (A) and Reason (R) are true, and R is the correct explanation of A
\item[\optlabel{B}] Both A and R are true, but R is \emph{not} the correct explanation of A
\item[\optlabel{C}] A is true, but R is false
\item[\optlabel{D}] A is false, but R is true
\end{itemize}
\end{qitem}
\qsep

\needspace{7\baselineskip}
\begin{qitem}{78}
\textbf{Statement I:} For $\varphi = \frac13$, phase estimation with any finite number $t$ of counting qubits never returns an
outcome $y$ with $y/2^t = \varphi$ exactly, with certainty or otherwise.\\[1mm]
\textbf{Statement II:} The number $\frac13$ has no terminating binary expansion ($\frac13 = 0.010101\ldots_2$).
\begin{itemize}[leftmargin=2.6em, labelsep=0.4em, itemsep=2pt, topsep=3pt]
\item[\optlabel{A}] Both Statement I and Statement II are true
\item[\optlabel{B}] Both Statement I and Statement II are false
\item[\optlabel{C}] Statement I is true, but Statement II is false
\item[\optlabel{D}] Statement I is false, but Statement II is true
\end{itemize}
\end{qitem}
\qsep

\needspace{20\baselineskip}
\begin{passagebox}{Paragraph for Questions 79 and 80}
\textbf{Shor's algorithm for $N = 21$ with $a = 2$.} A counting register of $t = 10$ qubits ($M = 2^{10} = 1024$) and a
5-qubit work register initialised to $\ket{1}$ are used. Steps: (i) $H^{\otimes 10}$ on the counting register;
(ii) $\ket{x}\ket{1} \mapsto \ket{x}\ket{2^x \bmod 21}$; (iii) inverse QFT on the counting register; (iv) measure the counting
register, obtaining $y$. Classical post-processing replaces $y/1024$ by the closest fraction $k/r$ with $r < 21$ (continued
fractions) and keeps $r$ if $2^r \equiv 1 \pmod{21}$. The true order of $2$ modulo $21$ is $6$, since $2^6 = 64 = 3\cdot21 + 1$.
\end{passagebox}
\needspace{7\baselineskip}
\begin{qitem}{79}
Which of the following measured values $y$ leads the post-processing to the correct order $r = 6$?
\par\smallskip\noindent\begin{tabularx}{\linewidth}{@{}*{4}{L}@{}}
\optlabel{A}\,$683$ & \optlabel{B}\,$512$ & \optlabel{C}\,$171$ & \optlabel{D}\,$341$
\end{tabularx}
\end{qitem}
\qsep

\needspace{7\baselineskip}
\begin{qitem}{80}
The probability that the measured value is exactly $y = 512$ is closest to
\par\smallskip\noindent\begin{tabularx}{\linewidth}{@{}*{4}{L}@{}}
\optlabel{A}\,$0.333$ & \optlabel{B}\,$0.500$ & \optlabel{C}\,$0.167$ & \optlabel{D}\,$0.001$
\end{tabularx}
\end{qitem}
\qsep

\newpage

\begin{center}{\Large\bfseries\color{brand} RESPONSE SHEET}\\[2pt]
{\small Darken exactly ONE bubble per question. Leave all four bubbles empty to skip a question (0 marks).}\end{center}
\textbf{Name:} \hrulefill\quad \textbf{Roll No.:} \rule{35mm}{0.4pt}\par\medskip
\noindent\begin{minipage}[t]{0.235\linewidth}\centering
\textbf{Section A}\\[2pt]
\begin{tabular}{@{}r@{\hspace{5pt}}l@{}}
\textbf{1} & \bub{A}\,\bub{B}\,\bub{C}\,\bub{D} \\[1.6pt]
\textbf{2} & \bub{A}\,\bub{B}\,\bub{C}\,\bub{D} \\[1.6pt]
\textbf{3} & \bub{A}\,\bub{B}\,\bub{C}\,\bub{D} \\[1.6pt]
\textbf{4} & \bub{A}\,\bub{B}\,\bub{C}\,\bub{D} \\[1.6pt]
\textbf{5} & \bub{A}\,\bub{B}\,\bub{C}\,\bub{D} \\[1.6pt]
\textbf{6} & \bub{A}\,\bub{B}\,\bub{C}\,\bub{D} \\[1.6pt]
\textbf{7} & \bub{A}\,\bub{B}\,\bub{C}\,\bub{D} \\[1.6pt]
\textbf{8} & \bub{A}\,\bub{B}\,\bub{C}\,\bub{D} \\[1.6pt]
\textbf{9} & \bub{A}\,\bub{B}\,\bub{C}\,\bub{D} \\[1.6pt]
\textbf{10} & \bub{A}\,\bub{B}\,\bub{C}\,\bub{D} \\[1.6pt]
\textbf{11} & \bub{A}\,\bub{B}\,\bub{C}\,\bub{D} \\[1.6pt]
\textbf{12} & \bub{A}\,\bub{B}\,\bub{C}\,\bub{D} \\[1.6pt]
\textbf{13} & \bub{A}\,\bub{B}\,\bub{C}\,\bub{D} \\[1.6pt]
\textbf{14} & \bub{A}\,\bub{B}\,\bub{C}\,\bub{D} \\[1.6pt]
\textbf{15} & \bub{A}\,\bub{B}\,\bub{C}\,\bub{D} \\[1.6pt]
\textbf{16} & \bub{A}\,\bub{B}\,\bub{C}\,\bub{D} \\[1.6pt]
\textbf{17} & \bub{A}\,\bub{B}\,\bub{C}\,\bub{D} \\[1.6pt]
\textbf{18} & \bub{A}\,\bub{B}\,\bub{C}\,\bub{D} \\[1.6pt]
\textbf{19} & \bub{A}\,\bub{B}\,\bub{C}\,\bub{D} \\[1.6pt]
\textbf{20} & \bub{A}\,\bub{B}\,\bub{C}\,\bub{D} \\[1.6pt]
\end{tabular}
\end{minipage}\hfill\begin{minipage}[t]{0.235\linewidth}\centering
\textbf{Section B}\\[2pt]
\begin{tabular}{@{}r@{\hspace{5pt}}l@{}}
\textbf{21} & \bub{A}\,\bub{B}\,\bub{C}\,\bub{D} \\[1.6pt]
\textbf{22} & \bub{A}\,\bub{B}\,\bub{C}\,\bub{D} \\[1.6pt]
\textbf{23} & \bub{A}\,\bub{B}\,\bub{C}\,\bub{D} \\[1.6pt]
\textbf{24} & \bub{A}\,\bub{B}\,\bub{C}\,\bub{D} \\[1.6pt]
\textbf{25} & \bub{A}\,\bub{B}\,\bub{C}\,\bub{D} \\[1.6pt]
\textbf{26} & \bub{A}\,\bub{B}\,\bub{C}\,\bub{D} \\[1.6pt]
\textbf{27} & \bub{A}\,\bub{B}\,\bub{C}\,\bub{D} \\[1.6pt]
\textbf{28} & \bub{A}\,\bub{B}\,\bub{C}\,\bub{D} \\[1.6pt]
\textbf{29} & \bub{A}\,\bub{B}\,\bub{C}\,\bub{D} \\[1.6pt]
\textbf{30} & \bub{A}\,\bub{B}\,\bub{C}\,\bub{D} \\[1.6pt]
\textbf{31} & \bub{A}\,\bub{B}\,\bub{C}\,\bub{D} \\[1.6pt]
\textbf{32} & \bub{A}\,\bub{B}\,\bub{C}\,\bub{D} \\[1.6pt]
\textbf{33} & \bub{A}\,\bub{B}\,\bub{C}\,\bub{D} \\[1.6pt]
\textbf{34} & \bub{A}\,\bub{B}\,\bub{C}\,\bub{D} \\[1.6pt]
\textbf{35} & \bub{A}\,\bub{B}\,\bub{C}\,\bub{D} \\[1.6pt]
\textbf{36} & \bub{A}\,\bub{B}\,\bub{C}\,\bub{D} \\[1.6pt]
\textbf{37} & \bub{A}\,\bub{B}\,\bub{C}\,\bub{D} \\[1.6pt]
\textbf{38} & \bub{A}\,\bub{B}\,\bub{C}\,\bub{D} \\[1.6pt]
\textbf{39} & \bub{A}\,\bub{B}\,\bub{C}\,\bub{D} \\[1.6pt]
\textbf{40} & \bub{A}\,\bub{B}\,\bub{C}\,\bub{D} \\[1.6pt]
\end{tabular}
\end{minipage}\hfill\begin{minipage}[t]{0.235\linewidth}\centering
\textbf{Section C}\\[2pt]
\begin{tabular}{@{}r@{\hspace{5pt}}l@{}}
\textbf{41} & \bub{A}\,\bub{B}\,\bub{C}\,\bub{D} \\[1.6pt]
\textbf{42} & \bub{A}\,\bub{B}\,\bub{C}\,\bub{D} \\[1.6pt]
\textbf{43} & \bub{A}\,\bub{B}\,\bub{C}\,\bub{D} \\[1.6pt]
\textbf{44} & \bub{A}\,\bub{B}\,\bub{C}\,\bub{D} \\[1.6pt]
\textbf{45} & \bub{A}\,\bub{B}\,\bub{C}\,\bub{D} \\[1.6pt]
\textbf{46} & \bub{A}\,\bub{B}\,\bub{C}\,\bub{D} \\[1.6pt]
\textbf{47} & \bub{A}\,\bub{B}\,\bub{C}\,\bub{D} \\[1.6pt]
\textbf{48} & \bub{A}\,\bub{B}\,\bub{C}\,\bub{D} \\[1.6pt]
\textbf{49} & \bub{A}\,\bub{B}\,\bub{C}\,\bub{D} \\[1.6pt]
\textbf{50} & \bub{A}\,\bub{B}\,\bub{C}\,\bub{D} \\[1.6pt]
\textbf{51} & \bub{A}\,\bub{B}\,\bub{C}\,\bub{D} \\[1.6pt]
\textbf{52} & \bub{A}\,\bub{B}\,\bub{C}\,\bub{D} \\[1.6pt]
\textbf{53} & \bub{A}\,\bub{B}\,\bub{C}\,\bub{D} \\[1.6pt]
\textbf{54} & \bub{A}\,\bub{B}\,\bub{C}\,\bub{D} \\[1.6pt]
\textbf{55} & \bub{A}\,\bub{B}\,\bub{C}\,\bub{D} \\[1.6pt]
\textbf{56} & \bub{A}\,\bub{B}\,\bub{C}\,\bub{D} \\[1.6pt]
\textbf{57} & \bub{A}\,\bub{B}\,\bub{C}\,\bub{D} \\[1.6pt]
\textbf{58} & \bub{A}\,\bub{B}\,\bub{C}\,\bub{D} \\[1.6pt]
\textbf{59} & \bub{A}\,\bub{B}\,\bub{C}\,\bub{D} \\[1.6pt]
\textbf{60} & \bub{A}\,\bub{B}\,\bub{C}\,\bub{D} \\[1.6pt]
\end{tabular}
\end{minipage}\hfill\begin{minipage}[t]{0.235\linewidth}\centering
\textbf{Section D}\\[2pt]
\begin{tabular}{@{}r@{\hspace{5pt}}l@{}}
\textbf{61} & \bub{A}\,\bub{B}\,\bub{C}\,\bub{D} \\[1.6pt]
\textbf{62} & \bub{A}\,\bub{B}\,\bub{C}\,\bub{D} \\[1.6pt]
\textbf{63} & \bub{A}\,\bub{B}\,\bub{C}\,\bub{D} \\[1.6pt]
\textbf{64} & \bub{A}\,\bub{B}\,\bub{C}\,\bub{D} \\[1.6pt]
\textbf{65} & \bub{A}\,\bub{B}\,\bub{C}\,\bub{D} \\[1.6pt]
\textbf{66} & \bub{A}\,\bub{B}\,\bub{C}\,\bub{D} \\[1.6pt]
\textbf{67} & \bub{A}\,\bub{B}\,\bub{C}\,\bub{D} \\[1.6pt]
\textbf{68} & \bub{A}\,\bub{B}\,\bub{C}\,\bub{D} \\[1.6pt]
\textbf{69} & \bub{A}\,\bub{B}\,\bub{C}\,\bub{D} \\[1.6pt]
\textbf{70} & \bub{A}\,\bub{B}\,\bub{C}\,\bub{D} \\[1.6pt]
\textbf{71} & \bub{A}\,\bub{B}\,\bub{C}\,\bub{D} \\[1.6pt]
\textbf{72} & \bub{A}\,\bub{B}\,\bub{C}\,\bub{D} \\[1.6pt]
\textbf{73} & \bub{A}\,\bub{B}\,\bub{C}\,\bub{D} \\[1.6pt]
\textbf{74} & \bub{A}\,\bub{B}\,\bub{C}\,\bub{D} \\[1.6pt]
\textbf{75} & \bub{A}\,\bub{B}\,\bub{C}\,\bub{D} \\[1.6pt]
\textbf{76} & \bub{A}\,\bub{B}\,\bub{C}\,\bub{D} \\[1.6pt]
\textbf{77} & \bub{A}\,\bub{B}\,\bub{C}\,\bub{D} \\[1.6pt]
\textbf{78} & \bub{A}\,\bub{B}\,\bub{C}\,\bub{D} \\[1.6pt]
\textbf{79} & \bub{A}\,\bub{B}\,\bub{C}\,\bub{D} \\[1.6pt]
\textbf{80} & \bub{A}\,\bub{B}\,\bub{C}\,\bub{D} \\[1.6pt]
\end{tabular}
\end{minipage}
\par\bigskip
\begin{center}\textbf{Score computation} (filled in by the evaluator, using the Answer Book)\\[4pt]
\renewcommand{\arraystretch}{1.5}
\begin{tabular}{|c|l|c|c|c|c|}\hline
\textbf{Sec.} & \textbf{Topic} & \textbf{Attempted} & \textbf{Correct} $C$ & \textbf{Wrong} $W$ & \textbf{Marks} $3C - W$ \\ \hline
A & Mathematics for the QFT & & & & \\ \hline
B & Quantum Mechanics & & & & \\ \hline
C & QFT Basics & & & & \\ \hline
D & Advanced QFT & & & & \\ \hline
\multicolumn{5}{|r|}{\textbf{Total (out of 240)}} & \\ \hline
\end{tabular}\end{center}

\end{document}